\documentclass[final]{beamer}

\usepackage[orientation=portrait,size=a0,scale=1.12]{beamerposter}

\usetheme[
  block=fill,
  sectionpage=none,
  progressbar=none,
  numbering=none,
]{metropolis}
\setbeamertemplate{navigation symbols}{}

\usepackage{iftex}
\ifPDFTeX
  \usepackage[T1]{fontenc}
  \usepackage{lmodern}
\fi

\usepackage{microtype}
\setlength{\emergencystretch}{3em}
\usepackage{amsmath,amssymb}
\usepackage{booktabs}
\usepackage{graphicx}
\usepackage{tikz}
\usetikzlibrary{arrows.meta,positioning,shapes.geometric}
\usepackage{qrcode}
\usepackage{xcolor}

\definecolor{PosterPrimary}{HTML}{0B3C5D}
\definecolor{PosterAccent}{HTML}{B31B1B}
\definecolor{PosterLight}{HTML}{F6F7F9}

\setbeamercolor{headline}{fg=white,bg=PosterPrimary}
\setbeamercolor{block title}{fg=white,bg=PosterPrimary}
\setbeamercolor{block body}{fg=black,bg=PosterLight}
\setbeamercolor{alertblock title}{fg=white,bg=PosterAccent}
\setbeamercolor{alertblock body}{fg=black,bg=PosterAccent!6}
\setbeamercolor{itemize item}{fg=PosterAccent}
\setbeamercolor{itemize subitem}{fg=PosterAccent}

\hypersetup{
  colorlinks=true,
  urlcolor=PosterAccent,
  linkcolor=PosterAccent,
  citecolor=PosterAccent,
}

\setbeamerfont{title}{size=\Huge,series=\bfseries}
\setbeamerfont{subtitle}{size=\Large}
\setbeamerfont{author}{size=\Large}
\setbeamerfont{institute}{size=\large}
\setbeamerfont{block title}{size=\Large,series=\bfseries}
\setbeamerfont{block body}{size=\normalsize}

\newlength{\sepwidth}
\newlength{\colwidth}
\setlength{\sepwidth}{0.02\textwidth}
\setlength{\colwidth}{0.32\textwidth}
\newcommand{\separatorcolumn}{\begin{column}{\sepwidth}\end{column}}
\setlength{\columnsep}{0pt}

\ifXeTeX
  \setsansfont[
    UprightFont={[FiraSans-Light.otf]},
    ItalicFont={[FiraSans-LightItalic.otf]},
    BoldFont={[FiraSans-Regular.otf]},
    BoldItalicFont={[FiraSans-Italic.otf]},
  ]{[FiraSans-Light.otf]}
  \setmonofont[
    UprightFont={[FiraMono-Regular.otf]},
    ItalicFont={[FiraMono-Oblique.otf]},
    BoldFont={[FiraMono-Medium.otf]},
    BoldItalicFont={[FiraMono-MediumOblique.otf]},
  ]{[FiraMono-Regular.otf]}
\fi

\newcommand{\posterincludegraphics}[2][]{%
  \IfFileExists{#2}{%
    \includegraphics[#1]{#2}%
  }{%
    \IfFileExists{figures/#2}{%
      \includegraphics[#1]{figures/#2}%
    }{%
      \IfFileExists{../presentation/figures/#2}{%
        \includegraphics[#1]{../presentation/figures/#2}%
      }{%
        \IfFileExists{../presentation/figures/crops/#2}{%
          \includegraphics[#1]{../presentation/figures/crops/#2}%
        }{%
          \IfFileExists{presentation/figures/#2}{%
            \includegraphics[#1]{presentation/figures/#2}%
          }{%
            \IfFileExists{presentation/figures/crops/#2}{%
              \includegraphics[#1]{presentation/figures/crops/#2}%
            }{%
              \IfFileExists{manuscript/images/#2}{%
                \includegraphics[#1]{manuscript/images/#2}%
              }{%
                \fbox{\parbox{0.95\linewidth}{\centering Missing image: \texttt{\detokenize{#2}}}}%
              }%
            }%
          }%
        }%
      }%
    }%
  }%
}

\newcommand{\fitimage}[1]{%
  \centering
  \posterincludegraphics[width=\linewidth,keepaspectratio]{#1}%
}

\newcommand{\posterpanel}[3]{%
  \begin{minipage}[t]{#1}
    \centering
    \posterincludegraphics[width=\linewidth]{#2}\\[0.2cm]
    {\small #3}
  \end{minipage}%
}

\title{Smoothly Coupled Remodeling}
\subtitle{Nonlocal stimulus \& fabric-guided anisotropy with robust numerics}
\author{Petr Heny\v{s}, Niels Hammer}
\institute{Technical University of Liberec (TUL)\\Medical University Graz (MUG)\\
{\normalsize \texttt{petr.henys@tul.cz}}}
\date{}

\setbeamertemplate{headline}{%
  \begin{beamercolorbox}[wd=\paperwidth,sep=1.2cm,left]{headline}
    \begin{columns}[T,totalwidth=\paperwidth]
      \begin{column}{0.71\paperwidth}
        {\usebeamerfont{title}\inserttitle\par}
        \vspace{0.2cm}
        {\usebeamerfont{subtitle}\insertsubtitle\par}
        \vspace{0.6cm}
        {\usebeamerfont{author}\insertauthor\par}
        {\usebeamerfont{institute}\insertinstitute\par}
      \end{column}
      \begin{column}{0.02\paperwidth}\end{column}
      \begin{column}{0.27\paperwidth}
        \raggedleft
        \qrcode[height=6.2cm]{https://doi.org/10.5281/zenodo.18242655}\par
        \vspace{0.4cm}
        {\large \href{https://doi.org/10.5281/zenodo.18242655}{doi:10.5281/zenodo.18242655}}\par
        \vspace{0.2cm}
        {\large \href{https://github.com/petr-henys/Smoothly-Coupled-Remodeling-with-Nonlocal-Stimulus-and-Fabric-Guided-Anisotropy}{GitHub repository}}\par
      \end{column}
    \end{columns}
  \end{beamercolorbox}
  \vspace{0.8cm}
}

\begin{document}
\begin{frame}[t]
  \begin{columns}[t,totalwidth=\textwidth]
    \begin{column}{\colwidth}
      \begin{block}{Motivation \& aim}
        \begin{itemize}\itemsep0.4em
          \item Stable, reproducible finite-element (FE) bone remodeling on \textbf{anatomical meshes}.
          \item \textbf{Three-field coupling:} density $\rho$, stimulus $S$, and log-fabric $\mathbf{L}=\ln\mathbf{A}$ (anisotropy).
          \item Key requirement: robustness to (i) \textbf{eigenvalue degeneracy} (orthotropy) and (ii) \textbf{stiff multiphysics coupling} (fixed-point solve).
        \end{itemize}
      \end{block}

      \begin{alertblock}{Coupled model (one implicit time step)}
        \vspace{-0.2cm}
        \begin{center}
          \resizebox{0.98\linewidth}{!}{%
            \begin{tikzpicture}[
              node distance=1.45cm,
              block/.style={draw, rounded corners, align=center, minimum width=7.2cm, minimum height=1.8cm, fill=white},
              small/.style={draw, rounded corners, align=center, minimum width=6.3cm, minimum height=1.6cm, fill=white},
              arrow/.style={-Latex, thick},
              every node/.style={transform shape},
            ]
              \node[block] (mech) {Mechanics\\ $-\nabla\!\cdot\!\sigma(\mathbf{u},\rho,\mathbf{L})=0$\\ $\Rightarrow\ \psi,\ \bar{\mathbf{Q}}$};
              \node[small, right=2.0cm of mech] (fab) {Fabric\\ $\partial_t\mathbf{L} - D_A\Delta\mathbf{L}$\\ $+\frac{a}{\tau_A}(\mathbf{L}-\mathbf{L}_M)=0$};
              \node[small, below=1.0cm of fab] (stim) {Stimulus\\ $\tau_S\partial_t S + S$\\ $-\nabla\!\cdot(\tau_S D_S\nabla S)=\mathrm{Drive}(\psi/\rho)$};
              \node[small, left=2.0cm of stim] (dens) {Density\\ $\partial_t\rho - \nabla\!\cdot(D_\rho\nabla\rho)$\\ $=\mathrm{Form}(S)-\mathrm{Res}(S)\rho$};

              \path[arrow] (mech) edge node[above, font=\large]{$\bar{\mathbf{Q}}$} (fab);
              \path[arrow] (mech) edge[out=-35, in=160] node[pos=0.45, right, font=\large]{$\psi$} (stim);
              \path[arrow] (stim) edge node[above, font=\large]{$S$} (dens);
              \path[arrow] (fab) edge[out=180, in=25] node[pos=0.35, left, font=\large]{$\mathbf{L}$} (mech);
              \path[arrow] (dens) edge[out=90, in=-155] node[pos=0.35, left, font=\large]{$\rho$} (mech);
            \end{tikzpicture}%
          }
        \end{center}
        \vspace{-0.2cm}
        \begin{itemize}\itemsep0.35em
          \item Backward Euler in time; FE discretization in space.
          \item Partitioned coupling solved by block Gauss--Seidel with optional Anderson mixing.
        \end{itemize}
      \end{alertblock}

      \begin{block}{What the simulation outputs}
        \begin{itemize}\itemsep0.35em
          \item \textbf{Density} $\rho(\mathbf{x},t)$: trabecular/cortical redistribution and CT-style comparisons.
          \item \textbf{Stimulus} $S(\mathbf{x},t)$: nonlocal mechanotransductive driver (diffusion--reaction length-scale).
          \item \textbf{Fabric} $\mathbf{A}(\mathbf{x},t)$ / $\mathbf{L}(\mathbf{x},t)$: anisotropy magnitude and principal directions.
          \item Diagnostics: fixed-point residuals, convergence histories, and sensitivity sweeps.
        \end{itemize}
      \end{block}

      \begin{block}{Reproducibility \& availability}
        \begin{itemize}\itemsep0.35em
          \item Code + scripts: \href{https://github.com/petr-henys/Smoothly-Coupled-Remodeling-with-Nonlocal-Stimulus-and-Fabric-Guided-Anisotropy.git}{GitHub repository}.
          \item Versioned archive (code, inputs, raw outputs): \href{https://doi.org/10.5281/zenodo.18242655}{Zenodo DOI}.
          \item Each run stores config snapshots, logs, and machine-readable metrics alongside field output.
        \end{itemize}
      \end{block}
    \end{column}

    \separatorcolumn

    \begin{column}{\colwidth}
      \begin{alertblock}{Contribution I: robust orthotropy across degeneracy}
        \begin{itemize}\itemsep0.35em
          \item \textbf{Problem:} near isotropy / transverse isotropy, eigenvectors are non-unique $\Rightarrow$ basis flips, discontinuous orthotropic response.
          \item \textbf{Approach:} evaluate spectral quantities via \textbf{Sylvester projectors} $\mathbf{P}_i$ (no eigenvectors), with smooth equalization in degenerate subspaces and enforced partition of unity.
          \item \textbf{Used for:} fabric-guided constitutive mapping and trace-free fabric target $\mathbf{L}_M(\bar{\mathbf{Q}})$.
        \end{itemize}
        \vspace{0.35cm}
        \fitimage{sylvester_projectors.png}
      \end{alertblock}

      \begin{alertblock}{Contribution II: robust coupling via safeguarded Anderson mixing}
        \begin{itemize}\itemsep0.35em
          \item Solve $\mathbf{x}=(\mathbf{L},S,\rho)$ by fixed-point iteration (block Gauss--Seidel).
          \item Stabilize stiff regimes with \textbf{Anderson (Pulay/DIIS) mixing} in scaled DOF space:
            Tikhonov regularization, conditioning- and stall-based restarts, step limiting, and fall-back to Picard when strongly contractive.
        \end{itemize}
        \vspace{0.35cm}
        \fitimage{anderson_diagnostic.png}
      \end{alertblock}

      \begin{block}{Verification (box benchmark)}
        \begin{itemize}\itemsep0.35em
          \item Refinement studies track field--norm pair errors for $(\mathbf{L},S,\rho)$ under mesh and time-step refinement.
          \item Regularization suppresses checkerboarding and improves mesh-insensitive trends.
        \end{itemize}
        \vspace{0.35cm}
        \fitimage{convergence_ab.png}
      \end{block}

      \begin{alertblock}{Take-home message}
        \begin{itemize}\itemsep0.35em
          \item Robust anatomical remodeling needs robustness in \textbf{spectral operations} (orthotropy) and \textbf{coupling iteration} (multiphysics).
          \item CT-template comparison on the proximal femur highlights both \textbf{set-point governance} (cortical/trabecular contrast) and \textbf{loading-direction sensitivity}.
          \item Sylvester-projector orthotropy + safeguarded Anderson mixing enable stable, reproducible simulations and sensitivity studies.
        \end{itemize}
      \end{alertblock}
    \end{column}

    \separatorcolumn

    \begin{column}{\colwidth}
      \begin{block}{Femur case study: CT template comparison}
        \begin{itemize}\itemsep0.35em
          \item Population CT-derived density template mapped to FE mesh (IDW) and rescaled to $[\rho_{\min},\rho_{\max}]$ (diagnostic similarity).
          \item Similarity metrics: RMSE and pattern correlation (also cortical/trabecular masks).
          \item Sensitivity axes: set-point ratio $r=\psi_{\mathrm{ref}}^{\mathrm{cort}}/\psi_{\mathrm{ref}}^{\mathrm{trab}}$ and hip force direction $\Delta\alpha_{\mathrm{front}}$.
        \end{itemize}
        \vspace{0.3cm}
        \begin{center}
          \posterpanel{\linewidth}{../results/psi_ref_ratio_sweep/psi_ref_ratio_diagnostic.png}{Set-point ratio sweep: CT agreement \& cortical--trabecular contrast}\par
          \vspace{0.55cm}
          \posterpanel{\linewidth}{../results/alpha_front_sweep/alpha_front_diagnostic.png}{Hip-direction sweep: CT agreement, pattern correlation, medial shift}\par
          \vspace{0.55cm}
          \posterpanel{0.49\linewidth}{alpha_front_rho.png}{Spatial density redistribution (angle sweep)}\hfill
          \posterpanel{0.49\linewidth}{../results/psi_ref_ratio_sweep/psi_ref_ratio_histograms.png}{Density distributions vs CT (final state)}
        \end{center}
      \end{block}
    \end{column}
  \end{columns}
\end{frame}
\end{document}
